\documentclass{beamer}
\usetheme{Madrid}
\usepackage{amsmath, amssymb, amsthm}
\usepackage{graphicx}
\usepackage{listings}
\usepackage{gensymb}
\usepackage[utf8]{inputenc}
\usepackage{hyperref}
\usepackage{gvv}
\usepackage{tikz}
\lstset{
  language=Python,
  basicstyle=\ttfamily\small,
  keywordstyle=\color{blue},
  stringstyle=\color{orange},
  numbers=left,
  numberstyle=\tiny\color{gray},
  breaklines=true,
  showstringspaces=false
}
\usetikzlibrary{decorations.pathmorphing}

\title{Question-2.10.51}
\author{EE25BTECH11022 - sankeerthan}
\date{}
\begin{document}

\frame{\titlepage}

\begin{frame}
\frametitle{Question}
If $\mathbf{a}, \mathbf{b}, \mathbf{c}$ are three non-zero, non-coplanar vectors and \\

$\mathbf{b}_1 = \mathbf{b} - \frac{\mathbf{b} \cdot \mathbf{a}}{|\mathbf{a}|^2} \mathbf{a}, \qquad \mathbf{b}_2 = \mathbf{b} + \frac{\mathbf{b} \cdot \mathbf{a}}{|\mathbf{a}|^2} \mathbf{a}$, \\
$\begin{aligned}
\mathbf{c}_1 = \mathbf{c} - \frac{\mathbf{c} \cdot \mathbf{a}}{|\mathbf{a}|^2} \mathbf{a} + \frac{\mathbf{b} \cdot \mathbf{c}}{|\mathbf{c}|^2} \mathbf{b}_1, 
\mathbf{c}_2 = \mathbf{c} - \frac{\mathbf{c} \cdot \mathbf{a}}{|\mathbf{a}|^2} \mathbf{a} + \frac{\mathbf{b}_1 \cdot \mathbf{c}}{|\mathbf{b}_1|^2} \mathbf{b}_1, \\
\mathbf{c}_3 = \mathbf{c} - \frac{\mathbf{c} \cdot \mathbf{a}}{|\mathbf{c}|^2} \mathbf{a} + \frac{\mathbf{b} \cdot \mathbf{c}}{|\mathbf{c}|^2} \mathbf{b}_1,
\mathbf{c}_4 = \mathbf{c} - \frac{\mathbf{c} \cdot \mathbf{a}}{|\mathbf{c}|^2} \mathbf{a} = \frac{\mathbf{b} \cdot \mathbf{c}}{|\mathbf{b}|^2} \mathbf{b}_1,
\end{aligned}$ \\

then the set of orthogonal vectors is
\begin{enumerate} 
\item($\mathbf{a}, \mathbf{b}_1, \mathbf{c}_3$) 
\item($\mathbf{a}, \mathbf{b}_1, \mathbf{c}_2$)
\item($\mathbf{a}, \mathbf{b}_1, \mathbf{c}_1$)
\item($\mathbf{a}, \mathbf{b}_2, \mathbf{c}_2$) 
\end{enumerate}
\end{frame}
\begin{frame}
\frametitle{Solution}
Let $\mathbf{a},\ \mathbf{b},\ \mathbf{c}$ be three linearly independent (non-coplanar) vectors in $\mathbb{R}^3$. 

We form a $3 \times 3$ matrix whose columns are these vectors:
\begin{align}
A = \myvec{| & | & | \\ \mathbf{a} & \mathbf{b} & \mathbf{c} \\ | & | & |}
\end{align}

To obtain an orthogonal set, apply the projection formula (Gram-Schmidt):

Let $\mathbf{v}_1 = \mathbf{a}$.
Orthogonalize $\mathbf{b}$ against $\mathbf{a}$:
\begin{align}
\mathbf{v}_2 = \mathbf{b}_1 = \mathbf{b} - \mathrm{proj}_{\mathbf{a}} \mathbf{b}
= \mathbf{b} - \frac{\mathbf{a}^\mathrm{T} \mathbf{b}}{\mathbf{a}^\mathrm{T} \mathbf{a}} \mathbf{a}
\end{align}
Orthogonalize $\mathbf{c}$ against $\mathbf{a}$ and $\mathbf{b}_1$:
\begin{align}
\mathbf{v}_3 = \mathbf{c}_1 
= \mathbf{c} - \frac{\mathbf{a}^\mathrm{T}\mathbf{c}}{\mathbf{a}^\mathrm{T}\mathbf{a}} \mathbf{a} 
  - \frac{\mathbf{b}_1^\mathrm{T}\mathbf{c}}{\mathbf{b}_1^\mathrm{T}\mathbf{b}_1} \mathbf{b}_1
\end{align}
\end{frame}
\begin{frame}
\frametitle{Solution}
Thus, the matrix with orthogonal columns is:
\begin{align}
Q = \myvec{| & | & | \\ \mathbf{a} & \mathbf{b}_1 & \mathbf{c}_1 \\ | & | & | }
\end{align}
The columns in $Q$ form an orthogonal set.The set of orthogonal vectors is:
($\mathbf{a},\mathbf{b}_1,\mathbf{c}_1$)
\end{frame}

\begin{frame}[fragile]
\frametitle{C-Code}
\begin{lstlisting}[language=C]
#include <stdio.h>
#include <math.h>

void dot_product(double *a, double *b, int dim, double *result) {
    *result = 0.0;
    for (int i = 0; i < dim; ++i) {
        *result += a[i] * b[i];
    }
}

void scalar_mult(double *a, double k, int dim, double *result) {
    for (int i = 0; i < dim; ++i) {
        result[i] = k * a[i];
    }
}

void vector_sub(double *a, double *b, int dim, double *result) {
    for (int i = 0; i < dim; ++i) {
        result[i] = a[i] - b[i];
    }
}
\end{lstlisting}
\end{frame}

\begin{frame}[fragile]
\frametitle{C-Code}
\begin{lstlisting}[language=C]
void gram_schmidt(double *a, double *b, double *c, double *b1, double *c3) {
    int dim = 3;
    double a_dot_b, a_dot_a, a_dot_c, c_dot_b1, b1_dot_b1;
    double proj_a_on_b[3], proj_a_on_c[3], proj_b1_on_c[3];

    dot_product(a, a, dim, &a_dot_a);
    dot_product(b, a, dim, &a_dot_b);
    scalar_mult(a, a_dot_b / a_dot_a, dim, proj_a_on_b);
    vector_sub(b, proj_a_on_b, dim, b1);

    dot_product(c, a, dim, &a_dot_c);
    scalar_mult(a, a_dot_c / a_dot_a, dim, proj_a_on_c);
    double temp[3];
    vector_sub(c, proj_a_on_c, dim, temp);

    dot_product(b1, b1, dim, &b1_dot_b1);
    dot_product(c, b1, dim, &c_dot_b1);
\end{lstlisting}
\end{frame}

\begin{frame}[fragile]
\frametitle{C-Code}
\begin{lstlisting}[language=C]
    scalar_mult(b1, c_dot_b1 / b1_dot_b1, dim, proj_b1_on_c);

    for (int i = 0; i < dim; ++i)
        c3[i] = temp[i] - proj_b1_on_c[i];}
__attribute__((visibility("default")))
void get_orthogonal_vectors(double *a, double *b, double *c, double *out_a, double *out_b1, double *out_c3) {
    for(int i=0; i<3; ++i) out_a[i]=a[i];
    gram_schmidt(a, b, c, out_b1, out_c3);
\end{lstlisting}
\end{frame}

\begin{frame}[fragile]
\frametitle{Python-Code}
\begin{lstlisting}
# Code by GVV Sharma
# Modified for Problem Solution
# Released under GNU GPL
# Calculating area enclosed between curves
# call.py

import ctypes
import numpy as np

lib = ctypes.CDLL('./code.so')
\end{lstlisting}
\end{frame}

\begin{frame}[fragile]
\frametitle{Python-Code}
\begin{lstlisting}
ib.get_orthogonal_vectors.argtypes = [
    ctypes.POINTER(ctypes.c_double),
    ctypes.POINTER(ctypes.c_double),
    ctypes.POINTER(ctypes.c_double),
    ctypes.POINTER(ctypes.c_double),
    ctypes.POINTER(ctypes.c_double),
    ctypes.POINTER(ctypes.c_double)
]

def get_orthogonal_vectors(a, b, c):
    a = np.array(a, dtype=np.float64)
    b = np.array(b, dtype=np.float64)
    c = np.array(c, dtype=np.float64)
\end{lstlisting}
\end{frame}
\begin{frame}[fragile]
\frametitle{Python-Code}
\begin{lstlisting}
 out_a = np.zeros(3, dtype=np.float64)
    out_b1 = np.zeros(3, dtype=np.float64)
    out_c3 = np.zeros(3, dtype=np.float64)
    lib.get_orthogonal_vectors(
        a.ctypes.data_as(ctypes.POINTER(ctypes.c_double)),
        b.ctypes.data_as(ctypes.POINTER(ctypes.c_double)),
        c.ctypes.data_as(ctypes.POINTER(ctypes.c_double)),
        out_a.ctypes.data_as(ctypes.POINTER(ctypes.c_double)),
        out_b1.ctypes.data_as(ctypes.POINTER(ctypes.c_double)),
        out_c3.ctypes.data_as(ctypes.POINTER(ctypes.c_double))
    )
    return out_a, out_b1, out_c3
\end{lstlisting}
\end{frame}

\end{document}